% !TeX root = ../../Book1.tex
% 纲要标题：### E.2 实闭域与实闭包
% 纲要位置：计划纲要.md，第 922 行。
% BEGIN APPENDIX OUTLINE SYNC
% #### E.2.1 实闭域的等价刻画
% #### E.2.2 实闭包的存在性与唯一性
% #### E.2.3 实闭性与序完备性
% END APPENDIX OUTLINE SYNC

\section{实闭域与实闭包}
\label{b1:sec:E02}

实数域并不代数闭：\(X^2+1\) 无实根。但在代数扩张中再加入任何缺少的根，便不能一直保留实数意义的正负比较。实闭性抽取这种代数上的极大性质，所需的结构定理与实数的序完备性必须分清。

% 纲要内容（仅作写作计划，尚非教材正文）：
%
% #### E.2.1 实闭域的等价刻画
% #### E.2.2 实闭包的存在性与唯一性
% #### E.2.3 实闭性与序完备性
%

\subsection{实闭域的等价刻画}
\label{b1:subsec:E02:01}

\begin{definition}\label{b1:def:real-closed-field}
若域 \(R\) 形式实，且没有真代数扩张仍为形式实域，则称 \(R\) 为\term[shibi-yu]{实闭域}[real closed field]。
\end{definition}

这个定义同时要求两件事：域本身可以赋序，而在它上面再作任何真代数扩张，都不可能赋予域序。限定“代数”很关键，因为这里讨论的是加入多项式的根，不是任意增加新的元素。例如 \(\mathbb Q\) 形式实，但 \(\mathbb Q(\sqrt2)\) 仍可以沿实数的序比较大小，所以 \(\mathbb Q\) 不实闭。“形式实”仅排除 \(-1\) 为平方和，“实闭”还要求这种性质不能在更大的代数域中继续保持。

\ProseParagraphBreak
以下定理把实闭域的极大性质改写成关于多项式的条件；后面的根计数证明将用到这些条件。等价刻画亦见\cite[Theorem A.5.19]{Mangolte2020RealAlgebraicVarieties}。
\begin{theorem}[实闭域的等价刻画]\label{b1:thm:real-closed-characterizations}
对形式实域 \(R\)，下列条件等价：
\begin{enumerate}
\item \(R\) 实闭；
\item \(R\) 有一个域序，使每个正元素都是平方，且 \(R[X]\) 中每个奇数次多项式都在 \(R\) 中有根；
\item 取 \(i^2=-1\)，二次扩张 \(R(i)\) 是代数闭域。
\end{enumerate}

这些条件成立时，\(R\) 的域序唯一，其正元素恰为非零平方。
\end{theorem}

\begin{proof}
先证第一项推出第二项。由\cref{b1:thm:artin-schreier-ordering}，给 \(R\) 取一个域序。若正元素 \(a\) 不是平方，则在二次扩张 \(R(\sqrt a)\) 中，等式 \(-1=\sum_j(x_j+y_j\sqrt a)^2\) 的常数项会给出
\[
-1=\sum_jx_j^2+a\sum_jy_j^2\ge0,
\]
矛盾。因此该扩张形式实，与实闭性矛盾，故每个正元素都是平方。

若有奇数次多项式无根，取次数大于一且最小的奇数次不可约多项式 \(f\)，记其次数为 \(n\)。域 \(R[X]/(f)\) 不是形式实，所以有次数均小于 \(n\) 的 \(g_j\)，使
\[
1+\sum_jg_j(X)^2=f(X)h(X).
\]
左侧的最高次系数为非零平方和，不会抵消，故其次数为偶数且至多 \(2n-2\)。于是 \(h\) 为奇数次多项式，次数小于 \(n\)。由 \(n\) 的最小性，\(h\) 的某个奇数次不可约因子必为一次，故 \(h\) 有根 \(b\in R\)。代入上式又得 \(1+\sum_jg_j(b)^2=0\)，矛盾。

再证第二项推出第三项。令 \(F=R(i)\)。\(F\) 中每个元素都是平方：对 \(a+bi\)，若 \(b\ne0\)，取正平方根 \(r=\sqrt{a^2+b^2}\)、\(u=\sqrt{(r+a)/2}\)，则 \(r>|a|\)、\(u>0\)，并有
\[
\left(u+\frac b{2u}i\right)^2=a+bi.
\]
若 \(b=0\)，按 \(a\) 的正负分别用实平方根或 \(i\sqrt{-a}\)，零元另取零。因而 \(F\) 没有二次扩张。

任取 \(F\) 上的一个多项式，令 \(N/R\) 为包含它的全部根及 \(F\) 的有限正规扩张。特征为零，故 \(N/R\) 为 Galois 扩张。取 \(G=\operatorname{Gal}(N/R)\) 的一个 Sylow 二子群 \(P\)。不动域 \(N^P/R\) 的次数为奇数；由\cref{b1:thm:primitive-element}，若这个扩张非平凡，就有次数大于一的奇数次不可约多项式，与第二项矛盾。因此 \(N^P=R\)、\(G=P\) 是二群。非平凡有限二群 \(H\) 总有指数二的子群：取极大真子群 \(M\)，\cref{b1:prop:p-normalizer-condition}使 \(M\trianglelefteq H\)，商若阶数大于二，Cauchy 定理给出它的非平凡真子群，与 \(M\) 极大矛盾。若 \(\operatorname{Gal}(N/F)\ne1\)，将此事实和 Galois 对应合用，便给出 \(F\) 的二次扩张，仍矛盾。故 \(N=F\)，原多项式的全部根已经在 \(F\) 中。这证明 \(F\) 代数闭。

第三项推出第一项时，若 \(L/R\) 为代数扩张，则\cref{b1:thm:embedding-extension}使它可嵌入代数闭域 \(R(i)\)。这个二次扩张没有其他中间域，所以 \(L\) 若为真扩张，便同构于 \(R(i)\)，包含 \(-1\) 的平方根而非形式实。因此 \(R\) 实闭。

最后，第二项说明正元素恰为非零平方，任何域序都必须把非零平方判为正，因此序唯一。
\end{proof}

三种条件分别适合不同的问题。第一项从扩张出发，说明不能再加入代数元而保持可赋序性；第二项允许直接检验两类方程，即正元素的平方根和奇数次多项式的根；第三项则说明只需加入一个 \(-1\) 的平方根，就足以使所有多项式分裂。第二项并未要求所有偶数次多项式有根，例如 \(X^2+1\) 始终无根，也未要求奇数次多项式只有一个根，例如 \(X^3-X\) 有三个根。

\ProseParagraphBreak
序唯一性意味着所有正负都由域运算决定。条件三也说明 \(R\) 自身并不代数闭，因为形式实性排除了 \(i\in R\)。

\begin{corollary}\label{b1:cor:real-polynomial-factorization}
实闭域 \(R\) 上的每个非零多项式都是非零常数、一次因子和形如 \((X-a)^2+b^2\)、\(a,b\in R\)、\(b\neq0\) 的二次因子的乘积。后面这些首一二次因子在 \(R\) 上处处为正。
\end{corollary}

\begin{proof}
由\cref{b1:thm:real-closed-characterizations}，多项式在 \(R(i)\) 中完全分裂。共轭映射 \(a+bi\mapsto a-bi\) 固定系数，所以非实根按共轭对出现，重数也相同。每对对应
\[
\bigl(X-(a+bi)\bigr)\bigl(X-(a-bi)\bigr)=(X-a)^2+b^2.
\]
实根给出一次因子。平方非负而 \(b^2>0\)，所以每个二次因子处处为正。
\end{proof}

因此实闭域上的不可约多项式只有一次和二次两类。没有实根的多项式未必不可约，例如 \(X^4+3X^2+2=(X^2+1)(X^2+2)\) 无实根，却分成两个二次因子。后面判断符号时真正使用的是：每个非实根及其共轭合成的首一二次因子永远为正，故只有一次因子能够让符号改变。

\subsection{实闭包的存在性与唯一性}
\label{b1:subsec:E02:02}

\begin{definition}\label{b1:def:real-closure}
设 \((K,\leq)\) 为有序域。若 \(R\) 是实闭域，\(K\hookrightarrow R\) 是保持序的域嵌入，且 \(R/K\) 为代数扩张，则称 \(R\) 连同该嵌入为 \((K,\leq)\) 的一个\term[shibi-bao]{实闭包}[real closure]。
\end{definition}

保持序的意思是：若 \(a,b\in K\) 且 \(a<b\)，它们在 \(R\) 中仍满足同一不等式。代数性则要求 \(R\) 的每个元素都满足一个系数在 \(K\) 中的非零多项式。寻找实闭包，就是在保留这些原有符号的前提下补足代数根，直到得到实闭域。不能直接使用代数闭包来代替，因为代数闭包含有 \(-1\) 的平方根，已经不能赋序。

\ProseParagraphBreak
下述定理固定的是有序域，不仅是遗忘了序的域。原始论证见\cite[Satz 8]{ArtinSchreier1927RealFields}。
\begin{theorem}[实闭包]\label{b1:thm:real-closure-existence}
每个有序域都有实闭包。若 \(R_1,R_2\) 是同一个有序域 \(K\) 的两个实闭包，则存在唯一保持序的 \(K\)-同构 \(R_1\rightarrow R_2\)。
\end{theorem}

\begin{proofsketch}
存在性可直接由极大有序代数扩张构造。固定 \(K\) 的代数闭包 \(\Omega\)，考虑其中包含 \(K\) 的子域及其延伸给定序的正锥，按子域包含和正锥延伸排序。链的子域与正锥分别取并，仍为有序代数扩张；Zorn 引理给出极大者 \(R\)。

核对极大者实闭，需要两项保持序的扩张事实。正元素 \(a\) 的平方根可以加入而保留序：否则由\cref{b1:thm:artin-schreier-ordering}证明末的正锥扩充方法，已有正锥生成的预序中会出现
\[
-1=\sum_jp_j(x_j+y_j\sqrt a)^2,\qquad p_j\ge0.
\]
比较常数项得到 \(-1=\sum_jp_j(x_j^2+ay_j^2)\ge0\)，矛盾。奇数次扩张同样保留序。若不然，取次数 \(n\) 最小的反例，用本原元写成 \(E[X]/(f)\)，其中 \(E\) 是所考虑的有序基域，\(n=\deg f\) 为奇数。不能扩充正锥意味着存在 \(p_j\ge0\)、\(\deg g_j<n\)，使
\[
1+\sum_jp_jg_j(X)^2=f(X)h(X).
\]
最高次项为正，故 \(\deg h\) 为小于 \(n\) 的奇数。取 \(h\) 的一个奇数次不可约因子及其根 \(\beta\)，代入上式，就在更小的奇数次扩张中得到同样的 \(-1\) 表示，与最小性矛盾。因此这两类扩张都能保留指定序。极大性使 \(R\) 的正元素已有平方根，奇数次多项式已有根；上一条等价刻画说明 \(R\) 实闭。

同构的存在性使用一个代数根的符号判定事实：给定有序域 \(E\)、不可约 \(f\in E[X]\) 以及它在实闭扩张中的第 \(j\) 个实根 \(\alpha\)，每个 \(g(\alpha)\)、\(g\in E[X]\)，的符号由 \(E\) 的序及这个根的序号决定。其证明对 \(f,g\) 和导数作带符号的 Euclid 余式计算，先数实根，再判定各根上的符号；遇到余式为零时降低次数继续，故只用有限次域运算和系数符号比较。这也是后一节 Sturm 方法的根符号判定推广。这里省略这项推广的余式表与各退化情形，实闭包的原始唯一性论证见\cite[Satz 8]{ArtinSchreier1927RealFields}，根符号判定和消去法见\cite[Sections 1.2, 1.4]{BochnakCosteRoy1998RealAG}。

现在取所有从 \(R_1\) 的中间域到 \(R_2\) 的保持序的部分 \(K\)-嵌入。链的并给出嵌入，取极大者 \(\phi:E\to R_2\)。若 \(\alpha\in R_1\setminus E\)，它的最小多项式经 \(\phi\) 变换后，在 \(R_2\) 中有同样多的实根；将 \(\alpha\) 送到相同序号的根。根符号判定使 \(g(\alpha)\) 的符号全部保持，因而这个域嵌入可保持序地延伸到 \(E(\alpha)\)，与极大性矛盾。故定义域为全部 \(R_1\)。像为实闭域，\(R_2\) 又在此像上代数且形式实，实闭性的极大性质使像等于 \(R_2\)。这样得到同构。

最后，保持序的 \(K\)-自同构会置换任意 \(f\in K[X]\) 的有限个实根，同时保持它们的次序，因此逐个固定。\(R_1/K\) 的代数性使它固定所有元素，故只能是恒等。把一个同构的逆与另一个复合，得到这样的自同构，因而同构唯一。
\end{proofsketch}

代数性不能删去。例如 \(\mathbb R\) 是包含 \(\mathbb Q\) 的实闭域，但不是 \(\mathbb Q\) 的实闭包。另一个容易混淆的条件是保持给定序：对 \(\mathbb Q(t)\)，取 \(t>0\) 的实闭包会使 \(t\) 成为平方，取 \(t<0\) 的实闭包则不可能如此，所以两者不能通过固定 \(t\) 的域同构相连。

\ProseParagraphBreak
唯一性中的“唯一同构”比一般代数闭包的唯一性更强，原因是保持序的同构不能互换一组有序的实根。

\subsection{实闭性与序完备性}
\label{b1:subsec:E02:03}

设 \((K,\leq)\) 为有序域。若 \(K\) 的每个非空且有上界的子集都在 \(K\) 中有上确界，即有一个最小上界，则称 \((K,\leq)\) 具有\term[xu-wanbeixing]{序完备性}[order completeness]。这个条件要求处理任意子集；实闭性则只规定有限个系数组成的多项式及其代数根。两个性质的要求不同，下面用具体域比较它们。

\ProseParagraphBreak
\(\mathbb R\) 的正元素有平方根，奇数次实多项式由通常的中间值定理有根，所以\cref{b1:thm:real-closed-characterizations}说明它实闭。此处只是利用实分析已经建立的性质来识别这个例子；下面从实闭性证明一般多项式的中间值性质，不以实数完备性作为一般有序域的公理。

\begin{proposition}\label{b1:prop:real-algebraic-closure}
实代数数域 \(\mathbb R_{\mathrm{alg}}=\{x\in\mathbb R\mid x\text{ 在 }\mathbb Q\text{ 上代数}\}\) 是 \(\mathbb Q\) 在通常序下的实闭包。
\end{proposition}

\begin{proof}
代数元在加、减、乘、除下封闭，所以这是域。其正元素的正平方根仍在 \(\mathbb R\) 中，并对原元素为二次代数，因此在 \(\mathbb Q\) 上代数。对一个系数在此域内的奇数次多项式，先在 \(\mathbb R\) 中取实根。有限个系数生成 \(\mathbb Q\) 的有限扩张，根对该有限扩张代数，由塔式公式在 \(\mathbb Q\) 上也代数。因此它满足\cref{b1:thm:real-closed-characterizations}的第二项，故实闭；它按定义在 \(\mathbb Q\) 上代数，且序与 \(\mathbb Q\) 相容，因而为实闭包。
\end{proof}

\(\mathbb R_{\mathrm{alg}}\) 没有实数域的序完备性。取任一实超越数 \(\tau\)，考虑子集 \(S=\{a\in\mathbb R_{\mathrm{alg}}\mid a<\tau\}\)。取比 \(\tau\) 小及大的整数，分别说明 \(S\) 非空并且有上界。若某个 \(s\in\mathbb R_{\mathrm{alg}}\) 是它的上确界，则 \(s\neq\tau\)。当 \(s<\tau\) 时，由有理数在实数中的稠密性，存在 \(q\in\mathbb Q\) 满足 \(s<q<\tau\)，于是 \(q\in S\)，说明 \(s\) 甚至不是上界。当 \(s>\tau\) 时，取 \(\tau<q<s\) 的有理数，便得到比 \(s\) 更小的上界。因此任何 \(s\) 都不可能是上确界。

\ProseParagraphBreak
设 \((K,\leq)\) 为有序域。若对每个 \(x\in K\)，都存在正整数 \(n\) 使 \(x<n\)，则称 \((K,\leq)\) 为\term[Archimedean-youxu-yu]{Archimedean 有序域}[Archimedean ordered field]。\(\mathbb Q\)、\(\mathbb R_{\mathrm{alg}}\) 与 \(\mathbb R\) 都满足这一性质。还可以取\secref{b1:sec:E01}中使 \(t>n\) 对所有正整数成立的 \(\mathbb Q(t)\) 的实闭包。由\cref{b1:thm:real-closure-existence}，这个实闭包存在且保留给定序，所以原不等式全部保留；其中的 \(t\) 仍大于每个正整数。因此实闭性既不蕴含 Archimedean 性，也不蕴含序完备性。
